\chapter{Reproduktion, Prüfprotokoll und offene Gates}
\index{Provenienz}\index{Release!Nachweise}
\section{Versionierte Auswertung}
Das Prüfwerkzeug verwendet ausschließlich die Standardbibliothek von Python:
\begin{quote}\path{scripts/verification_artifacts.py}\end{quote}
Es prüft
die 17 Nachweisgruppen, 72 historischen Messungen, Tabellen, Quellenhashes und
eingebetteten Abbildungen. Jede Änderung einer gebundenen Quelle erfordert
eine bewusste fachliche Prüfung und Neugenerierung; bloßes Aktualisieren eines
Hashes ist keine Abnahme. Die Originalzeiten werden nicht erneut gemessen.

\begin{lstlisting}[language=bash]
python scripts/verification_artifacts.py
python scripts/build_handbook.py --check-matrix
PYTHONPATH=src python -m pytest -q \
  tests/unit/test_verification_artifacts.py \
  tests/unit/test_uncertainty.py \
  tests/regression/test_uncertainty_study.py \
  tests/regression/test_accuracy_cost.py
\end{lstlisting}
Zur bewussten Neugenerierung werden die Entwicklungsabhängigkeiten und das
Visualisierungsextra benötigt. Der Befehl reproduziert die analytische
UQ-Stichprobe und zeichnet die eigenen Abbildungen:
\begin{lstlisting}[language=bash]
python scripts/verification_artifacts.py --write
PYTHONPATH=src python -m pytest -q
python scripts/build_handbook.py --self-test
\end{lstlisting}
Der letzte Befehl nutzt die festgelegte Docker/TeX-Werkzeugkette, baut alle fünf
PDFs und prüft negative Dokumentfälle. Ruff, Formatprüfung, mypy und Pakettests
ergänzen die numerischen Tests. Native PETSc-Tests brauchen die dokumentierte
native Umgebung; übersprungene Tests sind keine erfolgreichen Nachweise.

\section{Archiv und Herkunft}
\path{Doc/latex/verification-data.json} hält die vollständige historische
Umgebung, neu berechnete Gruppen/Budgetauswahlen, UQ-Ergebnisse und SHA-256
aller gebundenen Quellen sowie beider Abbildungen fest. Die bytegleiche
Benchmarkkopie \path{Doc/latex/evidence/m7-37c-linux-arm64.json} und der
synthetische UQ-Bericht werden mit den LaTeX-Quellen archiviert.
Die historischen Produktions- und Messcommits sind von dem auf der Titelseite
angegebenen Dokumentationscommit zu unterscheiden.

Das Handbucharchiv enthält die fertigen Tabellen und PDF-Abbildungen; zum
erneuten Setzen benötigt es keine Matplotlib-Installation. Die erneute
fachliche Auswertung benötigt dagegen das zugehörige Repository mit Quellen,
Tests und Auswertungsskript. Der Hash der UQ-Samplebytes ermöglicht den
Vergleich einer neu erzeugten Stichprobe; das Archiv behauptet keine
experimentellen Rohdaten für diesen Lehrfall.

\section{Tatsächlicher Reviewstatus}
Das Prüfprotokoll \path{Doc/latex/VV_VERIFICATION.md} dokumentiert lokale
Reproduktionen, Quellabgleich, Layoutprüfung und den zugeordneten CI-Stand.
Die neuen ausführbaren Prüfwerkzeuge folgen Red--Green--Refactor.
Die Dokumentation erhält keine künstlichen Tests, die nur ihre Formulierungen
duplizieren. Ein unabhängiges fachliches Review liegt nicht vor; eigene
Quellenprüfung und CI werden nicht als solches bezeichnet.

\section{Abnahme und Freigabe}
Dieser Band dokumentiert auch unerreichte Genauigkeitsziele, begrenzte
Backendprüfungen, unbekannte Fehlerbudgets und die fehlende experimentelle
Fallannahme. Sein vollständiger ehrlicher Abschluss kann Issue 210 erfüllen.
Die redaktionelle Gesamtabnahme aller Fachbände ist Issue 211; die Veröffentlichung
mit grünem Hauptzweig, Paketartefakten, Handbucharchiv, Prüfsummen und
annotiertem Versionstag liegt bei Issue 66. Die Titelseite bezeichnet bis dahin
den tatsächlichen Kandidatenstatus. Eine Dokumentationsabnahme allein ist keine
Veröffentlichung und keine experimentelle Validierung.
